%==============================================================================
% FICHAMENTO CRÍTICO — Automated Detection of Periodontal Bone Loss on
% Panoramic Radiographs Using Deep Learning
% Conforme NBR 6023 (referências) e NBR 10520 (citações)
%==============================================================================
\section{Fichamento: \citeonline{kurtbayrakdar2024periodontal}}

% --- 1. IDENTIFICAÇÃO ---
\subsection{Identificação}
\begin{tabular}{ll}
\textbf{Título:} & Detection of periodontal bone loss patterns and furcation
defects from panoramic radiographs using deep learning algorithm:
a retrospective study \\
\textbf{Autor(es):} & Kurt-Bayrakdar, Sevda; Bayrakdar, İbrahim Şevki;
Yavuz, Muhammet Burak; Sali, Nichal; Çelik, Özer; Köse, Oğuz; Uzun Saylan,
Bilge Cansu; Kuleli, Batuhan; Jagtap, Rohan; Orhan, Kaan \\
\textbf{Periódico:} & BMC Oral Health \\
\textbf{Ano:} & 2024 \\
\textbf{Volume:} & 24 \\
\textbf{Número:} & 1 \\
\textbf{Páginas:} & 155 \\
\textbf{DOI:} & \url{https://doi.org/10.1186/s12903-024-03896-5} \\
\textbf{Qualis:} & A2 (Odontologia) \\
\textbf{Tipo:} & Artigo original retrospectivo \\
\end{tabular}

% --- 2. OBJETIVO ---
\subsection{Objetivo}
Desenvolver e avaliar o desempenho de um algoritmo de deep learning baseado
em redes neurais convolutionais (CNN) com arquitetura U-Net para a detecção
automatizada de perda óssea alveolar total, padrões de perda óssea
interdental (horizontal e vertical) e defeitos de furca em radiografias
panorâmicas, visando criar um sistema de suporte ao diagnóstico
periodontal.

% --- 3. METODOLOGIA ---
\subsection{Metodologia}
\textbf{Desenho do estudo:} Estudo retrospectivo de desenvolvimento e
validação de algoritmo de deep learning com segmentação semântica. \\
\textbf{Amostra:} 1.121 radiografias panorâmicas com 2.251 marcações de
perda óssea alveolar total, 25.303 de perda óssea interdental (21.839
horizontal e 3.464 vertical) e 2.815 defeitos de furca. \\
\textbf{Métricas principais:} AUC, sensibilidade, precisão, F1-score,
acurácia, matriz de confusão. \\
\textbf{Validação:} Divisão treino/validação/teste; arquitetura U-Net
com 800 épocas; otimizador Adam; PyTorch; validação cruzada com
hiperparâmetros ajustados por parâmetro (learning rate 0,0001 para perda
óssea e horizontal; 0,00001 para vertical e furca).

% --- 4. RESULTADOS PRINCIPAIS ---
\subsection{Resultados Principais}
\begin{itemize}
  \item O sistema apresentou o maior desempenho na detecção de perda óssea
        alveolar total (AUC = 0,951; sensibilidade = 1,0; precisão = 0,995;
        acurácia = 0,994).
  \item Para perda óssea horizontal: AUC = 0,910; sensibilidade = 0,947;
        F1-score = 0,943; acurácia = 0,892.
  \item Para perda óssea vertical: AUC = 0,733; sensibilidade = 0,558;
        F1-score = 0,673 — o pior desempenho entre todas as classes.
  \item Para defeitos de furca: AUC = 0,868; sensibilidade = 0,892;
        precisão = 0,933; acurácia = 0,837.
\end{itemize}

% --- 5. LIMITAÇÕES ---
\subsection{Limitações}
\begin{itemize}
  \item O desempenho significativamente inferior na detecção de perda óssea
        vertical (AUC = 0,733) limita a aplicabilidade clínica do modelo
        para casos de defeitos infraósseos, que são clinicamente relevantes
        para o planejamento cirúrgico periodontal.
  \item Estudo retrospectivo unicêntrico, sem validação externa em
        populações com diferentes características demográficas e
        distribuições de doença periodontal.
  \item O padrão-ouro utilizado foi a anotação manual por observadores
        especialistas, introduzindo potencial viés de subjetividade na
        demarcação dos defeitos ósseos.
  \item Ausência de correlação com parâmetros clínicos periodontais
        (sondagem, nível de inserção clínica, sangramento) para confirmar
        os achados radiográficos.
\end{itemize}

% --- 6. CONTRIBUIÇÃO PARA O LIVRO ---
\subsection{Contribuição para o Livro}
Este artigo fundamenta diretamente o Capítulo 14, que aborda a aplicação
de deep learning em periodontia para detecção de perda óssea em
radiografias panorâmicas. O estudo é particularmente valioso por ser o
primeiro a utilizar segmentação semântica (U-Net) para distinguir padrões
de perda óssea (horizontal vs. vertical) e defeitos de furca, indo além
da simples classificação binária presença/ausência. O contraste entre o
alto desempenho para perda óssea total (AUC = 0,951) e o baixo desempenho
para perda vertical (AUC = 0,733) é discutido no livro como um desafio
aberto para o SDD, destacando a necessidade de arquiteturas específicas
para diferentes padrões de destruição periodontal.

% --- 7. ANÁLISE CRÍTICA ---
\subsection{Análise Crítica}
\begin{itemize}
  \item \textbf{Validade interna:} Alta -- Amostra robusta (1.121
        radiografias), arquitetura U-Net bem estabelecida, treinamento
        com 800 épocas e ajuste de learning rate por parâmetro.
  \item \textbf{Validade externa:} Média -- Dados de um único centro
        turco, sem validação em equipamentos de diferentes fabricantes
        ou populações etnicamente diversas.
  \item \textbf{Reprodutibilidade:} Alta -- Metodologia detalhada com
        especificação completa da arquitetura, hiperparâmetros e pipeline
        de treinamento em PyTorch.
  \item \textbf{Relevância clínica:} Alta -- A detecção automatizada de
        perda óssea total (AUC = 0,951) pode servir como ferramenta de
        triagem populacional; a discriminação horizontal/vertical auxilia
        no planejamento terapêutico.
  \item \textbf{Conflito de interesses:} Não -- Os autores declararam
        ausência de conflitos de interesse.
  \item \textbf{Viés identificado:} Viés de verificação parcial — os
        achados radiográficos não foram correlacionados com medidas
        clínicas periodontais (sondagem), que é o padrão-ouro para
        diagnóstico de periodontite.
\end{itemize}

% --- 8. CITAÇÕES RELEVANTES ---
\subsection{Citações Relevantes}
\begin{citacao}
``The system showed the highest diagnostic performance in the detection of
total alveolar bone losses (AUC = 0.951) and the lowest in the detection
of vertical bone losses (AUC = 0.733).'' (p. 6).
\end{citacao}

\begin{citacao}
``The sensitivity, precision, F1 score, accuracy, and AUC values were found
as 1, 0.995, 0.997, 0.994, 0.951 for total alveolar bone loss; ... 0.558,
0.846, 0.673, 0.506, 0.733 for vertical bone losses.'' (p. 7).
\end{citacao}

% --- 9. MAPEAMENTO SDD ---
\subsection{Mapeamento SDD}
\textbf{Problema:} Detecção e classificação automatizada de padrões de
perda óssea periodontal (total, horizontal, vertical) e defeitos de furca
em radiografias panorâmicas para auxílio ao diagnóstico periodontal. \\
\textbf{Entrada:} Radiografias panorâmicas anonimizadas (1.121 imagens)
com segmentação manual por especialistas como ground truth. \\
\textbf{Saída:} Mapas de segmentação semântica com identificação de perda
óssea total, horizontal, vertical e defeitos de furca por dente. \\
\textbf{Critérios:} AUC $\geq$ 0,85 para perda óssea total e horizontal;
AUC $\geq$ 0,75 para defeitos de furca; sensibilidade $\geq$ 0,55 para
perda vertical (aceitando desempenho inferior devido à complexidade do
padrão). \\
\textbf{Confiança:} 0,80 — Estudo com amostra ampla e metodologia robusta,
mas limitado pela natureza retrospectiva unicêntrica e ausência de
correlação clínica.
\end{tabular}

% --- 10. REFERÊNCIA ABNT ---
\subsection{Referência ABNT}
\begin{verbatim}
KURT-BAYRAKDAR, Sevda et al. Detection of periodontal bone loss patterns
and furcation defects from panoramic radiographs using deep learning
algorithm: a retrospective study. BMC Oral Health, London, v. 24, n. 1,
p. 155, 2024. DOI: 10.1186/s12903-024-03896-5.
\end{verbatim}
